\section{Fundamento te\'orico aplicado al PFC}
\subsection{2.1. Ingenier\'ia de Requisitos y necesidades}
La Ingenier\'ia de Requisitos (IR) permite transformar necesidades de personas y organizaciones en acuerdos que gu\'ien la construcci\'on de un sistema. No comienza redactando funciones por intuici\'on: identifica el contexto, distingue qu\'e problema es relevante y documenta requisitos que pueden revisarse. La norma ISO/IEC/IEEE 29148:2018 describe los procesos e informaci\'on necesarios para desarrollar requisitos de sistemas y software a lo largo del ciclo de vida \cite{iso29148}. El temario Foundation Level del IREB, versi\'on 3.1 indicada por la gu\'ia docente, complementa los conceptos de elicitaci\'on y especificaci\'on \cite{ireb2022}. Pohl presenta la IR como una pr\'actica que relaciona obtenci\'on, negociaci\'on, documentaci\'on, validaci\'on y gesti\'on en un marco coherente \cite{pohl2025}.

Para el PFC, pedir ``una aplicaci\'on que venda m\'as'' expresa una aspiraci\'on general, pero no indica qu\'e debe hacer la plataforma. Conviene primero examinar c\'omo el personal atiende una solicitud, consulta un producto y confirma su disponibilidad. A partir de esa informaci\'on pueden identificarse necesidades distintas: evitar registros duplicados, conocer existencias y revisar qu\'e recomendaciones son oportunas. Cada necesidad debe tener una fuente y quedar abierta a correcci\'on antes de convertirse en especificaci\'on.

\subsection{2.2. Niveles y tipos de requisitos}
Los requisitos de negocio representan resultados buscados por una organizaci\'on; los de interesados describen necesidades de grupos afectados; los de usuario permiten reconocer tareas; y los de sistema y software detallan las condiciones que el producto debe cumplir. La distinci\'on impide confundir una meta comercial con una decisi\'on tecnol\'ogica. Las gu\'ias de IR recomiendan mantener relaciones trazables entre niveles y no dar por aprobados supuestos sin responsable \cite{iso29148,swebok2024}.

\textbf{Aplicaci\'on:} una meta de negocio podr\'ia ser disminuir la repetici\'on de consultas de disponibilidad. Como necesidad del vendedor se propone ver en una sola pantalla la descripci\'on del art\'iculo, su precio y su stock. Un requisito funcional preliminar ser\'ia mostrar esos datos cuando el vendedor seleccione un c\'odigo de producto autorizado. Se trata de una propuesta acad\'emica: la meta, las pantallas y sus reglas de acceso requieren contraste con la empresa. No conviene definir desde el inicio una base de datos o una API concreta sin conocer las restricciones del PFC.

\subsection{2.3. Funciones, interfaces, restricciones y calidad}
Un requisito funcional indica un comportamiento comprobable. Por su parte, los requisitos de interfaz especifican intercambios con personas u otros sistemas, las restricciones acotan opciones de soluci\'on y los requisitos de calidad precisan las condiciones bajo las cuales el producto debe funcionar \cite{iso29148}. Glinz advierte que la etiqueta ``no funcional'' no tiene una interpretaci\'on universal y requiere una definici\'on cuidadosa en el proyecto \cite{glinz2007}. Por eso, en este informe se usa la etiqueta ``requisito de calidad'' cuando pueda asociarse con una caracter\'istica del modelo ISO/IEC 25010:2023 \cite{iso25010}.

Por ejemplo, ``el sistema ser\'a r\'apido'' no permite aceptar o rechazar el resultado. Un objetivo provisional m\'as preciso es que el 95\,\% de las consultas al cat\'alogo terminen en tres segundos o menos bajo una carga de ensayo acordada. El valor de tres segundos \emph{no es una medici\'on de ELECTRORIOS}; deber\'a negociarse y probarse. De forma semejante, un requisito de seguridad puede exigir que un vendedor no edite precios sin permiso, y uno de interfaz puede definir qu\'e datos se enviar\'ian al motor de recomendaci\'on sin incluir identificadores personales innecesarios. La necesidad de una interfaz real deber\'a corroborarse antes del dise\~no.

\subsection{2.4. Propiedades de buena formulaci\'on}
Un requisito individual debe ser necesario, apropiado, no ambiguo, completo, singular, factible, verificable, correcto y conforme al estilo acordado. Estas propiedades se discuten en la norma ISO/IEC/IEEE 29148:2018, especialmente en sus apartados 5.2.4--5.2.7 \cite{iso29148}. No basta con que una frase suene t\'ecnica: ``mostrar sugerencias adecuadas'' deja sin definir qu\'e es adecuado. Si el vendedor recibe una alternativa, se deber\'ia aclarar qu\'e informaci\'on acompa\~na la sugerencia y qu\'e acci\'on permite aceptarla o descartarla.

Un conjunto tambi\'en debe ser consistente, comprensible y susceptible de validaci\'on. Si un requisito permite vender sin existencias y otro lo impide, el conflicto debe resolverse con el responsable del proceso antes de programar. Las expresiones ``f\'acil'', ``siempre'' y ``lo m\'as pronto posible'' requieren condiciones medibles o deben reemplazarse. La plantilla EARS ayuda a representar un disparador y una respuesta observable sin mezclar varios comportamientos \cite{mavin2009}. Por ejemplo: ``Cuando el vendedor seleccione un producto activo, el sistema deber\'a mostrar el precio vigente registrado para ese producto''. Las decisiones sobre descuentos, impuestos o stock se a\~nadir\'an tras consultar la pol\'itica real del negocio.

\subsection{Relaci\'on con el marco y ciclo de vida}
Pohl propone considerar el contexto como parte central de la IR, y SWEBOK organiza el conocimiento del \`area de requisitos junto con otros campos de ingenier\'ia de software \cite{pohl2025,swebok2024}. La gu\'ia del docente tambi\'en solicita ISO/IEC/IEEE 15288:2023 para procesos de ciclo de vida de sistemas y la edici\'on ISO/IEC/IEEE 12207:2017 para software \cite{iso15288,iso12207}. Esta \`ultima se conserva aqu\'i como \emph{edici\'on solicitada en la gu\'ia}, aunque ISO ha publicado una versi\'on posterior. La selecci\'on definitiva de procesos se har\'a al conocer los acuerdos y recursos disponibles del PFC.

\noindent\textit{Nota de investigaci\'on:} las fichas normativas y los datos editoriales citados se verificaron en los portales de ISO, IEEE, Springer y IEEE Computer Society. Para una cita literal o interpretaci\'on m\'as espec\'ifica de una cl\'ausula, el equipo deber\'a contrastar la norma completa mediante el acceso institucional.
